\begin{figure}[p]
\centering
\newcommand{\lcol}{0.52\linewidth}
\newcommand{\rcol}{0.44\linewidth}

% ---------- Row 1 ----------
\begin{subfigure}[t]{\lcol}
\centering
\textbf{Quipper}
\begin{minted}{haskell}
bell :: Circ (Qubit,Qubit)
bell = do
  a <- qinit False; b <- qinit False
  hadamard a
  qnot_at b `controlled` a
  return (a,b)

with_computed bell $ \(a,b) -> do
  qnot_at a `controlled` b
  cdiscard (a,b)
\end{minted}
\end{subfigure}\hfill
\begin{subfigure}[t]{\rcol}
\centering
\textbf{Qunity}
\begin{minted}{text}
-- ASCII-only surface style
def deutsch(f : Bit -> Bit) : Bit =
  let x = plus in
  ctrl f(x) [
    0 -> x ;
    1 -> { gphase(pi); }
  ];
  had x;
  return x;
end
\end{minted}
\end{subfigure}

\vspace{2em}

% ---------- Row 2 (left only) ----------
\begin{subfigure}[t]{\linewidth}
\raggedright
\textbf{Q\#}
\begin{minted}{csharp}
operation CRy(theta : Double, ctrl : Qubit, tgt : Qubit)
  : Unit is Adj + Ctl {
    (Controlled Ry)([ctrl], (theta, tgt));
}

operation PrepareBell() : Unit {
    use (q0, q1) = (Qubit(), Qubit());
    H(q0);
    CNOT(q0, q1);
}
\end{minted}
\end{subfigure}

\vspace{2em}

% ---------- Row 3 ----------
\begin{subfigure}[t]{\lcol}
\centering
\textbf{OpenQASM 3}
\begin{minted}{c}
OPENQASM 3.0;
qubit q[2];

gate x(a) { U(pi, 0, pi) a; }
gate h(a) { U(pi/2, 0, pi) a; gphase(-pi/4); }

h q[0];
ctrl @ x q[0], q[1];   // CNOT via modifier
\end{minted}

\end{subfigure}\hfill
\begin{subfigure}[t]{\rcol}
\centering\textbf{Stim}
\begin{minted}{text}
H 0
H 1
MPP X0*X1
DETECTOR rec[-1]
OBSERVABLE_INCLUDE(0) rec[-1]
\end{minted}
\end{subfigure}

\vspace{2em}


% ---------- Full-width bottom ----------
\begin{subfigure}[t]{\linewidth}
\raggedright
\textbf{QIR}
\begin{minted}{llvm}
%q0 = call %Qubit* @__quantum__rt__qubit_allocate()
%q1 = call %Qubit* @__quantum__rt__qubit_allocate()

call void @__quantum__qis__h__body(%Qubit* %q0)
call void @__quantum__qis__cnot__body(%Qubit* %q0, %Qubit* %q1)

call void @__quantum__rt__qubit_release(%Qubit* %q1)
call void @__quantum__rt__qubit_release(%Qubit* %q0)
\end{minted}
\end{subfigure}

\caption{
Six minimal snippets foregrounding each approach to representing circuits:
(a) Quipper’s boxed subcircuits and controlled combinators;
(b) Qunity’s coherent control over quantum data;
(c) Q\# functors for controlled/adjoint variants;
(d) OpenQASM~3 gate definitions and modifiers;
(e) Stim’s stabilizer format with parity measurements and detectors;
(f) QIR’s LLVM calls to quantum intrinsics.
}
\label{fig:lang-contrast-6}
\end{figure}
